\documentclass[12pt]{article}

% Layout.
\usepackage[top=1in, bottom=0.75in, left=1in, right=1in, headheight=1in, headsep=6pt]{geometry}

% Fonts.
\usepackage{mathptmx}
\usepackage[scaled=0.86]{helvet}
\renewcommand{\emph}[1]{\textsf{\textbf{#1}}}

% TiKZ.
\usepackage{tikz, pgfplots}
\usetikzlibrary{calc}

% Misc packages.
\usepackage{amsmath,amssymb,latexsym}
\usepackage{graphicx}
\usepackage{array}
\usepackage{xcolor}
\usepackage{multicol}
\usepackage{fancyhdr}
\pagestyle{fancy}


% Misc.
\renewcommand{\d}{\displaystyle}
\newcommand{\ds}{\displaystyle}
\def\bc{\begin{center}}
\def\ec{\end{center}}
\def\be{\begin{enumerate}}
\def\ee{\end{enumerate}}

\newcommand{\ans}[1][1in]{\rule{#1}{.5pt}}


\lhead{Math F113X: Quiz 3}

\begin{document}

\vspace{1cm}

\noindent \textbf{Name:} \hrulefill \quad  score:\ans[1cm] \ / 10 \\


There are 10 points possible on this quiz. No aids (book, notes, etc.)
are permitted. You may use a non-programmable calculator.  {\bf Show
all work and supporting calculations for full credit. Explain how you
get your answers.}



\be
\item Alden and Bobby buy a large pizza together for
\$32. It is split equally into one-third veggie, one-third anchovy,
and one-third pineapple.
\be
\item (2 points) Alden doesn't like anchovies, so he values the anchovy part at \$0, and he values pineapple and veggie
  equally. Bobby values the veggie \emph{twice} as much as the anchovy and
  he values the pineapple \emph{five times} as much as the anchovy. Complete the
  following table to indicate the value each party would assign to each
  part of the pizza.

\begin{center}
  \begin{tabular}{c||c|c|c||c}
    party & veggie & anchovy & pineapple & total value \\ \hline \hline
    Alden & & & & \\ \hline
    Bobby & & & &
  \end{tabular}
\end{center}

\vspace{2cm}

\item (3 points) Suppose they split the pizza using the \emph{Divider-Chooser}
  method
  with Alden as the divider. He divides the pizza into Share 1, which
  consists of the anchovy and veggie portions of the pizza and Share 2,
  which consists of the pineapple portion of the pizza. Complete the
  Divider-Chooser method, including the final outcome for each
  party. Be sure to fully justify your answer and verify that each party
  ends up with a fair share.
\ee

\vfill

Alden's final outcome \ans \ Bobby's final outcome \ans
\newpage

\item Linda, Ned, and Sophie will split a bag of jelly beans worth \$30
  using the {\bf Lone Divider} method. One of these people is selected
  to divide the jelly beans into three shares. The values are shown
  below.
  \begin{center}
    \begin{tabular}{c||c|c|c}
      Person & Share 1 & Share 2 & Share 3 \\ \hline
      Linda & \$8 & \$15 & \$7 \\
      Ned & \$10 & \$10 & \$10 \\
      Sophie & \$5 & \$20 & \$5
    \end{tabular}
  \end{center}

  \be
\item (1 point) What is the value of a fair share to each person?
  \vfill
  
\item (1 point) Who was the divider? How can you tell?
  \vfill
  
\item (1 points) \emph{Circle the bids (dollar values) in the table} that are
  fair shares to each person.
  \vspace{0.5cm}
  
\item (2 points) State the final allocation of the shares or describe what
  happens next in the Lone Divider process.
  \vfill
  \vfill
  \ee
  
\ee
\end{document}
%%% Local Variables:
%%% mode: LaTeX
%%% TeX-master: t
%%% End:
